% Preamble for CMF research papers.
% Defines the document class, fonts, theorem styles, and unicode mappings.
% Used by paper1.tex / paper2.tex / paper3.tex (which inline this content).
% To regenerate a self-contained paperN.tex from a body file:
%   cat preamble.tex paperN.body.tex > paperN.tex
% Compile with: xelatex paperN.tex
%
% This preamble does NOT include the \title/\author commands —
% those are at the top of each paperN.body.tex (paper-specific).

\documentclass[11pt]{article}

\usepackage[margin=1in]{geometry}
\usepackage{amsmath,amssymb,amsthm}
\usepackage{mathtools}
\usepackage{bm}

% Use xelatex with TeX Gyre Termes (Times equivalent) for body and math
\usepackage{fontspec}
\usepackage{unicode-math}
\setmainfont{TeX Gyre Termes}
\setmonofont{DejaVu Sans Mono}[Scale=0.85]
\setmathfont{Latin Modern Math}[Scale=MatchLowercase]

\usepackage{textcomp}
\usepackage{booktabs}
\usepackage{array}
\usepackage{longtable}
\usepackage{calc}
\usepackage{etoolbox}
\usepackage{makecell}
\usepackage[hidelinks]{hyperref}
\usepackage{enumitem}
\usepackage{xcolor}
\usepackage{framed}
\usepackage{fvextra}

% Enable pandoc's syntax-highlighted code-block environment.

% Tighter spacing
\setlength{\parskip}{0.5\baselineskip}
\setlength{\parindent}{0pt}

% Section heading style: bold, smaller than default. The markdown text already
% includes the section number ("## 1. Introduction"), so we suppress LaTeX's
% auto-numbering here.
\usepackage{titlesec}
\titleformat{\section}[block]{\normalsize\bfseries}{}{0pt}{}
\titleformat{\subsection}[block]{\normalsize\bfseries}{}{0pt}{}
\titleformat{\subsubsection}[block]{\normalsize\bfseries}{}{0pt}{}
\titlespacing*{\section}{0pt}{1.2\baselineskip}{0.4\baselineskip}
\titlespacing*{\subsection}{0pt}{1.0\baselineskip}{0.3\baselineskip}

% Pandoc-required commands
\providecommand{\tightlist}{\setlength{\itemsep}{0pt}\setlength{\parskip}{0pt}}
\providecommand{\pandocbounded}[1]{#1}

% Custom title block (centered, matching v2 style)
\makeatletter
\renewcommand\maketitle{%
  \begingroup
  \begin{center}
    \LARGE\bfseries \@title \par
  \end{center}
  \ifx\@subtitle\@empty\else
    \vspace{0.4em}
    \begin{center}
      \itshape \@subtitle \par
    \end{center}
  \fi
  \vspace{0.6em}
  \begin{center}
    \@author \par
    \vspace{0.2em}
    Correspondence: \@correspondence \par
    \vspace{0.2em}
    \@date \par
  \end{center}
  \vspace{1em}
  \endgroup
}
\def\@subtitle{}
\newcommand{\subtitleX}[1]{\def\@subtitle{#1}}
\def\@correspondence{}
\newcommand{\correspondenceX}[1]{\def\@correspondence{#1}}
\makeatother

% Custom abstract: bold heading centered, indented body
\renewenvironment{abstract}{%
  \begin{center}\bfseries Abstract\end{center}%
  \begin{quote}\setlength{\parindent}{0pt}}{\end{quote}\medskip}

% Keywords command
\newcommand{\keywordsX}[1]{%
  \begin{quote}\setlength{\parindent}{0pt}\textbf{Keywords:} #1\end{quote}}

\begin{document}
% Paper-specific metadata
\title{An Investigation of an Alternate Base for a Conservative Matrix
Field and an Example of a Search for a Representation of ζ(5)}
\subtitleX{Research Note --- Exploratory Mathematics}
\author{Louis Blaine}
\correspondenceX{louisblaine45@gmail.com}
\date{March 2026}

% Body
\maketitle

\begin{abstract}
Conservative Matrix Fields (CMFs), recently introduced by Raz et
al.~(2025) as a unifying structure underlying hundreds of formulas for
mathematical constants such as \(\pi\), are defined on a two-dimensional
integer lattice with a commutativity condition on mixed lattice paths.
In this note we investigate what happens when the underlying base is
extended from the real integer lattice \(\mathbb{Z}^{2}\) to pairs of
complex numbers over the Gaussian integers \(\mathbb{Z}\){[}i{]}, and
ask whether the resulting algebraic structure generates new mathematics.
We show that the 4-periodicity of i under multiplication naturally
stratifies lattice paths into two classes: commutative paths (4-step
closure) and anticommutative paths (2-step closure), the latter being
algebraically equivalent to spinors. The combined integrability
conditions admit a natural realisation in the Pauli algebra, yielding a
structure analogous to a superfield. We identify known constant pairs
(even/odd zeta values, alternating/non-alternating series) as
projections of this double-coverage structure, and derive predictions
about an irrationality proof for \(\zeta\)(5). A computational search
for an Apéry-like degree-5 recurrence guided by the anticommutative CMF
framework is described and partially executed, consistent with the
structural predictions while identifying the remaining computational
challenge.
\end{abstract}

\keywordsX{Conservative Matrix Field, Riemann zeta function,
irrationality, Clifford algebra, spinors, Apéry's theorem, continued
fractions, Ramanujan Machine}

\hypertarget{introduction}{%
\section{1. Introduction}\label{introduction}}

In late 2025, Raz, Shalyt, Leibtag, Kalisch, Weinbaum, Hadad, and
Kaminer {[}1{]} announced a remarkable discovery: a previously unknown
mathematical structure, the Conservative Matrix Field (CMF), underlies
hundreds of formulas for the constant \(\pi\), including those of
Archimedes, Euler, Madhava, and Ramanujan. Their Euler2AI project,
combining large language models with novel symbolic algorithms,
extracted 385 unique \(\pi\)-formulas from 455,050 arXiv preprints and
demonstrated that 43\% of them belong to a single CMF. This work grew
from the Ramanujan Machine project {[}2{]}, which uses AI to generate
conjectures about fundamental constants in the form of continued
fractions.

A CMF is defined on the two-dimensional integer lattice
\(\mathbb{Z}^{2}\). At each lattice point (m,n) one assigns 2\(\times\)2
step matrices M\_x(m,n) and M\_y(m,n) whose entries are polynomial
functions of m and n.\textsuperscript{1} The conservativity condition
--- the discrete analog of curl = 0 --- demands that mixed lattice steps
commute:\textsuperscript{2}

\[M_y(m+1,n) \cdot M_x(m,n) = M_x(m,n+1) \cdot M_y(m,n)\]

The Ramanujan Machine group showed that by choosing polynomial functions
f(x,y) and f*(x,y) defining the step matrices, one obtains CMFs whose
continued-fraction limits are fundamental constants. This note asks:
what happens if we replace the real integer lattice with a complex base?
Specifically, we investigate the Gaussian integer lattice
\(\mathbb{Z}\){[}i{]} and show that the algebraic structure of i --- its
4-periodicity under multiplication --- naturally produces two distinct
classes of lattice paths with fundamentally different integrability
conditions. We call these the commutative and anticommutative CMF
structures, and we argue that they correspond respectively to the even
and odd special values of the Riemann zeta function.

\hypertarget{the-periodicity-structure-of-i-and-path-classification}{%
\section{2. The Periodicity Structure of i and Path
Classification}\label{the-periodicity-structure-of-i-and-path-classification}}

The imaginary unit i satisfies the 4-cycle i\textsuperscript{1}=i,
i\textsuperscript{2}=−1, i\textsuperscript{3}=−i,
i\textsuperscript{4}=1. When lattice paths are taken over
\(\mathbb{Z}\){[}i{]} rather than \(\mathbb{Z}^{2}\), closed loops
naturally acquire different periodicity.

\textbf{Definition 2.1} (Path Classification). Let \(\gamma\) be an
elementary square loop in the lattice. We say \(\gamma\) is a
\emph{commutative path} if the accumulated matrix product around
\(\gamma\) equals the identity after 4 steps, and an
\emph{anticommutative path} if it equals minus the identity after 2
steps.

Setting S(m,n) = M\_y(m+1,n)\(\cdot\)M\_x(m,n) and T(m,n) =
M\_x(m,n+1)\(\cdot\)M\_y(m,n), the two integrability conditions are:

\[S - T = 0 \quad \text{(commutative CMF)}; \qquad S + T = 0 \quad \text{(anticommutative CMF)}\]

A field satisfying both conditions simultaneously requires {[}M\_x,
M\_y{]} = 0 and \{M\_x, M\_y\} = 0, equivalently M\_xM\_y = M\_yM\_x =
0. The Pauli matrices \(\sigma_{1}\), \(\sigma_{3}\) provide a natural
nontrivial solution: writing M\_x = a\(\sigma_{1}\) + b\(\sigma_{3}\),
M\_y = c\(\sigma_{2}\), both conditions are satisfied by the Clifford
relations \(\sigma _{i}\sigma _{j}\) + \(\sigma _{j}\sigma _{i}\) =
2\(\delta _{ij}\)I. Other solutions exist (e.g.~pairs of orthogonal
projectors such as diag(1,0) and diag(0,1)), but the Pauli realisation
is the one compatible with the nonzero polynomial step-matrix structure
required by the CMF construction.

\hypertarget{explicit-step-matrices-worked-example-for-c}{%
\subsection{\texorpdfstring{2.2. Explicit step matrices: worked example
for
c\textsubscript{1}}{2.2. Explicit step matrices: worked example for c}}\label{explicit-step-matrices-worked-example-for-c}}

We give the explicit step matrices for the constant c\textsubscript{1}
(see companion paper {[}3{]}, §3.1). Define:

\[M_x(m,n) = m\cdot\sigma_z + \sigma_x = \begin{bmatrix} m & 1 \\ 1 & -m \end{bmatrix} \quad \text{(trace-zero, antidiagonal-plus-diagonal)}\]

\[M_y(m,n) = \sigma_y\text{-scaled} = \begin{bmatrix} 0 & 1 \\ -1 & 0 \end{bmatrix} \quad (= i\cdot\sigma_y \text{ up to phase})\]

\textbf{Anticommutative CMF condition.} Compute S + T at any lattice
point (m,n):

\[S = M_y(m+1,n)\cdot M_x(m,n) = \begin{bmatrix}0&1\\-1&0\end{bmatrix}\begin{bmatrix}m&1\\1&-m\end{bmatrix} = \begin{bmatrix}1&-m\\-m&-1\end{bmatrix}\]

\[T = M_x(m,n+1)\cdot M_y(m,n) = \begin{bmatrix}m&1\\1&-m\end{bmatrix}\begin{bmatrix}0&1\\-1&0\end{bmatrix} = \begin{bmatrix}-1&m\\m&1\end{bmatrix}\]

\[S + T = \begin{bmatrix}0&0\\0&0\end{bmatrix}. \quad \checkmark^4\]

\textbf{Scalar recursion.} The projective (Möbius) action of M\_x(k,0)
on \(\mathbb{C}\mathbb{P}^{1}\) is z \(\mapsto\) (kz+1)/(z−k), so the
path product along the x-axis \(\prod\)\emph{\{k=1\}\^{}\{N\} M\_x(k,0)
generates the recursion z\_k = (kz}\{k-1\}+1)/(z\_\{k-1\}−k), i.e.~p(k)
= 1, r(k) = k in the notation of Theorem 2.1 of {[}3{]}. This is the
recurrence whose angle-sum limit gives S\textsubscript{1} =
arg(\(\Gamma\)(1+i/2)/\(\Gamma\)(1/2+i/2)) and c\textsubscript{1} =
\textbar tan(S\textsubscript{1})\textbar{} \(\approx\) 0.55500\ldots{}

\begin{longtable}[]{@{}lll@{}}
\toprule\noalign{}
k & M\_x(k,0) & Möbius step z \(\to\) (kz+1)/(z−k) \\
\midrule\noalign{}
\endhead
\bottomrule\noalign{}
\endlastfoot
1 & {[}{[}1,1{]},{[}1,−1{]}{]} & z \(\to\) (z+1)/(z−1) \\
2 & {[}{[}2,1{]},{[}1,−2{]}{]} & z \(\to\) (2z+1)/(z−2) \\
3 & {[}{[}3,1{]},{[}1,−3{]}{]} & z \(\to\) (3z+1)/(z−3) \\
k & {[}{[}k,1{]},{[}1,−k{]}{]} & z \(\to\) (kz+1)/(z−k) \\
\end{longtable}

\textbf{Table 0.} M\_x(k,0) for the first values of k, giving the
c\textsubscript{1} recurrence.

The corresponding M\_y matrices (all equal to
{[}{[}0,1{]},{[}−1,0{]}{]}) play no role in the path along the x-axis
but are required for the CMF integrability condition to hold on the full
2D lattice. Analogous constructions yield p(k)/r(k) = (2k−1)/k for
c\textsubscript{4} and k(k+1)/k\textsuperscript{2} for
c\textsubscript{5}; see {[}3, §3.3--3.4{]}.

\hypertarget{conventions}{%
\subsection{2.3. Conventions}\label{conventions}}

We use the principal branch Arg : \(\mathbb{C}\)* \(\to\) (−\(\pi\),
\(\pi\){]} throughout. The CMF constant is defined as c =
\textbar tan(Arg(z))\textbar{} = \textbar Im(z)/Re(z)\textbar{} where z
= \(\Gamma\)(1+w/2)/\(\Gamma\)(1/2+w/2) (the second form is manifestly
branch-free and equal to the first when Re(z) \(\neq\) 0). Numerical
evidence confirms Re(z) \textgreater{} 0 for all Gaussian primes of norm
\(\leq\) 41 (verified at 120-digit precision; see companion paper {[}4,
§2.3{]}). The parameter w = \(\pi\)/N(\(\pi\)) \(\in\) \(\mathbb{Q}\)(i)
for every Gaussian prime \(\pi\).

\hypertarget{double-coverage-real-and-complex-computational-paths}{%
\section{3. Double Coverage: Real and Complex Computational
Paths}\label{double-coverage-real-and-complex-computational-paths}}

The two path types create a double coverage of the constant landscape.
Along a commutative path (real lattice steps only), matrix products
accumulate as real polynomial products and converge to real constants
with all-positive-term series. Along an anticommutative path (steps
through the i-direction), each pair of steps introduces a sign flip via
i\textsuperscript{2} = −1, generating an alternating
series.\textsuperscript{5} This provides a structural explanation for
why certain pairs of mathematical constants are in fact projections of a
single complex CMF:

\begin{longtable}[]{@{}
  >{\raggedright\arraybackslash}p{(\columnwidth - 4\tabcolsep) * \real{0.3333}}
  >{\raggedright\arraybackslash}p{(\columnwidth - 4\tabcolsep) * \real{0.3333}}
  >{\raggedright\arraybackslash}p{(\columnwidth - 4\tabcolsep) * \real{0.3333}}@{}}
\toprule\noalign{}
\begin{minipage}[b]{\linewidth}\raggedright
Commutative partner
\end{minipage} & \begin{minipage}[b]{\linewidth}\raggedright
Anticommutative partner
\end{minipage} & \begin{minipage}[b]{\linewidth}\raggedright
Relationship
\end{minipage} \\
\midrule\noalign{}
\endhead
\bottomrule\noalign{}
\endlastfoot
\(\zeta\)(2) = \(\pi^{2}\)/6 & \(\pi\)/4 = 1−1/3+1/5−⋯ & \(\pi^{2}\) vs
\(\pi\) (degree drop) \\
\(\zeta\)(2n) --- closed form in \(\pi\) & \(\zeta\)(2n−1) --- no closed
form & Even vs odd \\
Li\textsubscript{2}(1) = \(\pi^{2}\)/6 & Li\textsubscript{2}(1/2) =
\(\pi^{2}\)/12 − ln\textsuperscript{2}(2)/2 & Full vs half path \\
e = \(\Sigma\) 1/n! & 1/e = \(\Sigma\) (−1)\textsuperscript{n}/n! &
Forward vs reverse \\
\end{longtable}

\textbf{Table 1:} Known constant pairs as commutative/anticommutative
CMF partners.

The degree-drop pattern (\(\pi^{2}\) paired with \(\pi\)) is
structurally significant: commutative paths close after 4 steps
(accumulating squared quantities), while anticommutative paths close
after 2 steps (accumulating linear quantities). This correctly predicts
that even zeta values \(\zeta\)(2n) are paired with odd values
\(\zeta\)(2n−1). The even zeta values all have closed forms in powers of
\(\pi\): \(\zeta\)(2)=\(\pi^{2}\)/6, \(\zeta\)(4)=\(\pi^{4}\)/90,
\(\zeta\)(6)=\(\pi^{6}\)/945. The odd zeta values
\(\zeta\)(3),\(\zeta\)(5),\(\zeta\)(7),\ldots{} have no such forms ---
not accidentally, but because they live on anticommutative paths.

\hypertarget{degree-5-anticommutative-cmf-and-predictions-for-5}{%
\section{\texorpdfstring{4. Degree-5 Anticommutative CMF and Predictions
for
\(\zeta\)(5)}{4. Degree-5 Anticommutative CMF and Predictions for (5)}}\label{degree-5-anticommutative-cmf-and-predictions-for-5}}

Apéry's 1978 proof {[}4{]} of the irrationality of \(\zeta\)(3) relied
on the three-term recurrence:

\[n^3 u_n = (34n^3 - 51n^2 + 27n - 5)\, u_{n-1} - (n-1)^3\, u_{n-2}\]

whose ratio b\_n/a\_n converges to \(\zeta\)(3) fast enough to prove
irrationality. Under the anticommutative CMF framework, we derive the
following predictions for \(\zeta\)(5):

\textbf{Prediction 4.1.} The irrationality proof for \(\zeta\)(5), if it
exists via this method, will use a degree-5 recurrence
n\textsuperscript{5} u\_n = A\textsubscript{5}(n) u\_\{n-1\} −
(n−1)\textsuperscript{5} u\_\{n-2\} where \(\alpha\) = lim\_n
A\textsubscript{5}(n)/n\textsuperscript{5} \textgreater{}
e\textsuperscript{5} \(\approx\) 148.41.

\textbf{Prediction 4.2.} The natural form is A\textsubscript{5}(n) =
(2n−1){[}Q\(\cdot n^{2}(n−1)^{2}\) + P\(\cdot\)n(n−1) + R{]} with
\(\alpha\) = 2Q, requiring Q \(\geq\) 75.

\textbf{Prediction 4.3.} (Poincaré-Perron) The characteristic equation
x\textsuperscript{2} − \(\alpha\)x + 1 = 0 must have roots \(\lambda\)+
\textgreater{} e\textsuperscript{5} \textgreater{} 1 \textgreater{}
\(\lambda\)− \textgreater{} 0 with \(\lambda+\lambda\)− = 1, giving
\textbar b\_n/a\_n − \(\zeta\)(5)\textbar{} \textasciitilde{}
(1/\(\lambda\)+)\textsuperscript{n} sufficient to prove irrationality
after clearing lcm(1,\ldots,n)\textsuperscript{5}.\textsuperscript{3}

\textbf{Prediction 4.4.} (Heuristic obstruction) Apéry's proof of
\(\zeta\)(3) succeeded because \(\zeta\)(4) = \(\pi^{4}\)/90 \(\neq\) 0
provided a non-vanishing commutative anchor. For \(\zeta\)(5), the
functional equation maps s=5 to \(\zeta\)(−4)=0 (a trivial zero),
removing this anchor. This heuristic obstruction, within our framework,
suggests why \(\zeta\)(5) is harder than \(\zeta\)(3), though it does
not constitute a proof that no Apéry-type recurrence exists.

\hypertarget{computational-search}{%
\section{5. Computational Search}\label{computational-search}}

We implemented a computational search for the predicted recurrence,
first validating it on Apéry's \(\zeta\)(3) recurrence: iterating to
n=50, the ratio b\textsubscript{50}/a\textsubscript{50} matches
\(\zeta\)(3) to 55 significant figures.

\hypertarget{search-design}{%
\subsection{5.1. Search Design}\label{search-design}}

Parameters Q \(\in\) {[}75, 119{]}, P \(\in\) {[}−300, 300{]}, R \(\in\)
{[}1, 60{]}; approximately 2.5 million candidates evaluated in 90
seconds. No recurrence satisfying all integrability and irrationality
conditions was found. The closest candidate was Q=75, P=−222, R=27,
which converges to within 8.75\(\times\)10\textsuperscript{-8} of
\(\zeta\)(5) but fails the integrality condition
(lcm(1,\ldots,n)\textsuperscript{5}\(\cdot\)a\_n \(\notin\)
\(\mathbb{Z}\) for n \(\geq\) 4).\textsuperscript{6}

\begin{longtable}[]{@{}
  >{\raggedright\arraybackslash}p{(\columnwidth - 4\tabcolsep) * \real{0.3333}}
  >{\raggedright\arraybackslash}p{(\columnwidth - 4\tabcolsep) * \real{0.3333}}
  >{\raggedright\arraybackslash}p{(\columnwidth - 4\tabcolsep) * \real{0.3333}}@{}}
\toprule\noalign{}
\begin{minipage}[b]{\linewidth}\raggedright
Property
\end{minipage} & \begin{minipage}[b]{\linewidth}\raggedright
\(\zeta\)(3) --- Apéry
\end{minipage} & \begin{minipage}[b]{\linewidth}\raggedright
\(\zeta\)(5) --- Predicted
\end{minipage} \\
\midrule\noalign{}
\endhead
\bottomrule\noalign{}
\endlastfoot
Recurrence degree & 3 & 5 \\
Leading coeff \(\alpha\) & 34 & 2Q \textgreater{} 148.41 \\
\(\lambda\)+ & 17+12\(\sqrt{2}\) \(\approx\) 33.97 & \textgreater{}
e\textsuperscript{5} \(\approx\) 148.41 \\
lcm exponent & 3 & 5 \\
Free params in A & 0 (fully determined) & 3 (Q, P, R) \\
Commutative anchor & \(\zeta\)(4)=\(\pi^{4}\)/90 \(\neq\) 0 &
\(\zeta\)(6)=\(\pi^{6}\)/945, func. eq. \(\to\) \(\zeta\)(−4)=0 \\
Status & Proven irrational (1978) & Open problem \\
\end{longtable}

\textbf{Table 2:} Comparison of Apéry's recurrence for \(\zeta\)(3) and
predicted structure for \(\zeta\)(5).

\hypertarget{scope-of-remaining-search}{%
\subsection{5.2. Scope of remaining
search}\label{scope-of-remaining-search}}

The feasible band has \(\lambda\)+ just above e\textsuperscript{5};
extending the search to Q \(\leq\) 1000 would require approximately 35
minutes on a 100-core cluster and is recommended as the immediate next
computational step.

\hypertarget{relation-to-the-linear-forms-approach-to-odd-zeta-values}{%
\subsection{5.3. Relation to the linear-forms approach to odd zeta
values}\label{relation-to-the-linear-forms-approach-to-odd-zeta-values}}

The negative search result at Q \(\leq\) 119 should be read against the
linear-forms approach to odd zeta values. Rivoal {[}10{]} established
that infinitely many odd zeta values \(\zeta\)(2n+1) are irrational, and
Zudilin {[}11{]} refined this to the statement that at least one of
\(\zeta\)(5), \(\zeta\)(7), \(\zeta\)(9), \(\zeta\)(11) is irrational.
Ball and Rivoal {[}12{]} proved the lower bound dim\_\(\mathbb{Q}\)
span\_\(\mathbb{Q}\)\{1, \(\zeta\)(3), \(\zeta\)(5), \ldots,
\(\zeta\)(2n+1)\} \(\geq\) (1 − o(1)) \(\cdot\) (log n)/(1 + log 2), and
the direction has since been sharpened by Fischler and others {[}13{]}.
These results establish irrationality of \(\mathbb{Q}\)-linear
combinations of odd zeta values via small linear forms, and do not
directly address the existence of an Apéry-type three-term recurrence
for any individual value --- the target of the CMF search presented
here. A negative outcome at our parameter range is therefore consistent
with the Rivoal--Zudilin framework and leaves open the existence of a
single-value recurrence for \(\zeta\)(5) at larger parameters.

\hypertarget{a-superfield-like-structure-and-targets-for-new-constants}{%
\section{6. A Superfield-like Structure and Targets for New
Constants}\label{a-superfield-like-structure-and-targets-for-new-constants}}

A matrix field with both commutative and anticommutative components has
the algebraic structure analogous to a superfield (in the sense of a
\(\mathbb{Z}_{2}\)-graded decomposition into even and odd sectors). In
the Pauli basis:

\[M_{\text{total}} = \zeta(6)^{1/2}\sigma_3 + \zeta(5)\sigma_1 + \Xi(5)\sigma_2\]

where \(\Xi\)(5) is the imaginary Pauli (\(\sigma_{2}\)) component
carrying the cross-term between sectors. If confirmed, \(\Xi\)(5) would
be a new fundamental constant accessible only through the
anticommutative CMF structure and invisible to existing commutative
methods. We do not formalize the \(\mathbb{Z}_{2}\)-graded structure as
a rigorous supersymmetry algebra here; the analogy is structural rather
than physical.

\hypertarget{missed-formulas-and-the-anticommutative-search}{%
\section{\texorpdfstring{7. Missed \(\pi\) Formulas and the
Anticommutative
Search}{7. Missed  Formulas and the Anticommutative Search}}\label{missed-formulas-and-the-anticommutative-search}}

The Euler2AI project {[}1{]} found that 43\% of known \(\pi\) formulas
belong to a single commutative CMF. We conjecture that a significant
fraction of the remaining 57\% belong to anticommutative CMFs ---
formulas involving alternating series, half-integer arguments, or
odd-weight modular forms invisible to a purely commutative search.
Running the Euler2AI pipeline with the anticommutative integrability
condition is the natural complement to the existing project.

\hypertarget{discussion-and-open-questions}{%
\section{8. Discussion and Open
Questions}\label{discussion-and-open-questions}}

\textbf{(1) Existence of the \(\zeta\)(5) recurrence.} Does a degree-5
recurrence satisfying all conditions exist? The framework suggests yes,
but parameter space is not yet fully explored.

\textbf{(2) Uniqueness.} Poincaré-Perron analysis suggests the system is
exactly determined (5 equations in 5 unknowns), pointing toward
uniqueness once a valid recurrence is found.

\textbf{(3) Higher odd zeta values.} The pattern predicts degree-7 for
\(\zeta\)(7), degree-9 for \(\zeta\)(9), with computational difficulty
growing as e\^{}d.

\textbf{(4) Mock theta functions.} Ramanujan's mock theta functions
carry half-integer modular weight --- a spinorial character --- and may
fit naturally within the anticommutative CMF framework.

\textbf{(5) Clifford generalisation.} Higher-dimensional Clifford
algebras might produce CMF structures connecting mathematical constants
to physical quantities. Whether the collapse mechanism from matrix
dynamics to scalar recursion (§2.2) extends to Cl(p,q) beyond Cl(2) is
an open question not addressed here.

\hypertarget{conclusion}{%
\section{9. Conclusion}\label{conclusion}}

We have shown that extending the base of a Conservative Matrix Field
from the real integer lattice to the Gaussian integers naturally
produces two classes of lattice paths --- commutative and
anticommutative --- whose integrability conditions admit a natural
realisation in the Pauli algebra, yielding a spinor-like algebraic
structure. Explicit step matrices are constructed (§2.2) and the
anticommutative CMF condition S+T=0 is verified. The commutative and
anticommutative projections correspond respectively to even and odd
special values of the Riemann zeta function. A computational search
covering \textasciitilde2.5 million candidates found no valid
\(\zeta\)(5) recurrence within the explored range, consistent with the
prediction that the true recurrence, if it exists, lies at larger
parameter values requiring cluster computation. The negative result is
consistent with the Rivoal--Zudilin framework on irrationality of
combinations of odd zeta values {[}10--12{]}, which targets a different
object (linear combinations) than the single-value Apéry-type recurrence
sought here.

\textbf{Computational reproducibility.} The recurrence search code and
Lean 4 formalization are available at
github.com/Loublain/cmf-gaussian-primes {[}archived:
doi:10.5281/zenodo.19670531{]}. The Lean formalization covers: CMF
structure and path independence (\texttt{CMF.lean},
\texttt{PathIndep.lean}), polynomial CMF (\texttt{PolyCMF.lean}),
commutative/anticommutative structure (\texttt{SuperCMF.lean}), Pauli
step matrix verification (\texttt{PauliCMF.lean}), zeta pairing
algebraic structure (\texttt{ZetaPairs.lean}), and recurrence
predictions (\texttt{Recurrence.lean}). All theorems are
machine-verified with no \texttt{sorry} and depend only on standard
Mathlib axioms (\texttt{propext}, \texttt{Classical.choice},
\texttt{Quot.sound}). All searches were implemented in Python with
mpmath (mp.dps = 50 for convergence checks) and are fully reproducible
from the repository README.

\hypertarget{references}{%
\section{References}\label{references}}

{[}1{]} T. Raz, M. Shalyt, E. Leibtag, R. Kalisch, S. Weinbaum, Y.
Hadad, I. Kaminer. From Euler to AI: Unifying Formulas for Mathematical
Constants. NeurIPS 2025; arXiv:2502.17533 (2025).

{[}2{]} G. Raayoni, S. Gottlieb, Y. Manor, et al.~Generating conjectures
on fundamental constants with the Ramanujan Machine. \emph{Nature} 590,
67--73 (2021). doi:10.1038/s41586-021-03229-4

{[}3{]} L. Blaine. On the Tangent of the Argument: Closed Forms for
Constants Arising from Anticommutative Pauli CMF Matrix Products.
Preprint (2026).

{[}4{]} R. Apéry. Irrationalité de \(\zeta\)(2) et \(\zeta\)(3).
\emph{Astérisque} 61, 11--13 (1979).

{[}5{]} O. David. The conservative matrix field. arXiv:2303.09318
(2023).

{[}6{]} G. Rhin, C. Viola. The group structure for \(\zeta\)(3).
\emph{Acta Arithmetica} 97, 269--293 (2001). doi:10.4064/aa97-3-6

{[}7{]} E. Witten. Supersymmetry and Morse theory. \emph{Journal of
Differential Geometry} 17(4), 661--692 (1982).

{[}8{]} M. Atiyah. On the K-theory of compact Lie groups.
\emph{Topology} 4(1), 95--99 (1965). doi:10.1016/0040-9383(65)90051-4

{[}9{]} S. Ramanujan. The Lost Notebook and Other Unpublished Papers.
Narosa Publishing House, New Delhi (1988).

{[}10{]} T. Rivoal. La fonction zêta de Riemann prend une infinité de
valeurs irrationnelles aux entiers impairs. \emph{C. R. Acad. Sci.
Paris} 331 (2000), 267--270. doi:10.1016/S0764-4442(00)01624-4

{[}11{]} W. Zudilin. One of the numbers \(\zeta\)(5), \(\zeta\)(7),
\(\zeta\)(9), \(\zeta\)(11) is irrational. \emph{Russian Math. Surveys}
56 (2001), 774--776. doi:10.1070/RM2001v056n04ABEH000427

{[}12{]} K. Ball, T. Rivoal. Irrationalité d'une infinité de valeurs de
la fonction zêta aux entiers impairs. \emph{Invent. Math.} 146 (2001),
193--207. doi:10.1007/s002220100168

{[}13{]} S. Fischler. Irrationalité de valeurs de zêta. \emph{Séminaire
Bourbaki}, Astérisque 294 (2004), 27--62.

\textbf{Acknowledgements.} The author thanks Claude Sonnet 4.6
(Anthropic) for assistance with derivation, step-matrix construction,
computational search, Lean 4 formalization, and manuscript preparation.
Correspondence: louisblaine45@gmail.com

\hypertarget{footnotes}{%
\section{Footnotes}\label{footnotes}}

\textsuperscript{1} The precise definition uses
\texttt{MvPolynomial\ (Fin\ 2)\ R} evaluated via
\texttt{Matrix.map\ (MvPolynomial.eval\ v)} where
\texttt{v\ :\ Fin\ 2\ →\ R} maps the two variables to the lattice
coordinates. The Lean 4 formalization is the canonical reference:
github.com/Loublain/cmf-gaussian-primes

\textsuperscript{2} Formally:
\texttt{My\ (p\ +\ (1,0))\ *\ Mx\ p\ =\ Mx\ (p\ +\ (0,1))\ *\ My\ p} for
all \texttt{p\ :\ ℤ\ ×\ ℤ}. Machine-verified in \texttt{CMF.lean}.

\textsuperscript{3} Prediction 4.3 is formally proved as
\texttt{Prediction43\_holds} in \texttt{Recurrence.lean}. Predictions
4.1 and 4.2 remain open conjectures. The Lean 4 formalization is
available at github.com/Loublain/cmf-gaussian-primes

\textsuperscript{4} The anticommutative condition S + T = 0 for the
c\textsubscript{1} step matrices is machine-verified in
\texttt{PauliCMF.lean} as \texttt{S\_plus\_T\_c1}.

\textsuperscript{5} The sign flip S = −T is formally proved as
\texttt{AntiCMF.sign\_flip} in \texttt{ZetaPairs.lean}. The connection
between this algebraic mechanism and convergence to specific zeta values
is not formalized; it requires analytic infrastructure not currently in
Mathlib.

\textsuperscript{6} The parameter space and closest candidate are
formally defined in \texttt{Recurrence.lean} as \texttt{SearchSpace} and
\texttt{closest\_candidate}, with membership verified by
\texttt{closest\_in\_search\_space}.

\end{document}
